% -*- coding: utf-8 -*-

\begin{zhaiyao}
\begin{spacing}{1.5}
{
水下场景重建面临光学退化与动态变化的双重挑战。现有静态水下重建方法能够显式处理水体退化，但难以建模动态扰动；通用动态新视角合成方法虽可表示时变几何，却通常忽略水下成像物理，容易产生颜色失真和高斯漂浮伪影。针对上述问题，本文以4DGaussians为动态场景表示骨架，将SeaSplat中的水下图像形成模型迁移至动态高斯渲染流程，构建面向扰动水下环境的新视角合成模型。模型引入了后向散射网络和直接透射衰减网络，用于估计水体介质对颜色的影响；同时引入Depth Anything V2生成的单目逆深度先验，约束高斯体空间分布。实验在Robot、Fish、Coral和Streaks四个动态水下场景以及SeaThru-NeRF静态水下数据集上进行。动态场景实验中，本文方法在Robot、Coral和Streaks三个场景上取得最优PSNR，平均PSNR达到30.54 dB，优于UDR-GS的30.05 dB和4DGS的29.63 dB；在SeaThru-NeRF数据集上与4DGS基线基本持平或更优。无参考图像质量评价（NIQE）结果表明，本文方法在多数场景的去水重建图像上取得了最优统计自然度，进一步验证了水下物理模型与深度监督的有效性。消融实验发现，水下物理建模与单目深度监督的联合引入并非在所有场景中均带来正向收益，二者存在一定的优化干扰，完整配置并不总是最优。
}
\end{spacing}
\end{zhaiyao}


\begin{guanjianci}
三维高斯溅射；水下场景重建；动态场景建模；水下物理成像；深度监督
\end{guanjianci}


\begin{abstract}
\begin{spacing}{1.5}
Underwater scene reconstruction faces the dual challenges of optical degradation and dynamic variation. Existing static underwater reconstruction methods can explicitly model water-induced degradation, but they struggle to represent dynamic disturbances. General dynamic novel view synthesis methods can represent time-varying geometry, but they usually ignore underwater imaging physics, which may lead to color distortion and floating Gaussian artifacts. To address these issues, this thesis adopts 4DGaussians as the dynamic scene representation backbone and transfers the underwater image formation model in SeaSplat into the dynamic Gaussian rendering pipeline, constructing a novel view synthesis model for disturbed underwater environments. The model introduces a backscatter network and a direct transmission attenuation network to estimate the influence of the water medium on color. Meanwhile, monocular inverse-depth priors generated by Depth Anything V2 are incorporated to constrain the spatial distribution of Gaussian primitives. Experiments are conducted on four dynamic underwater scenes, namely Robot, Fish, Coral, and Streaks, as well as the static SeaThru-NeRF underwater dataset. In the dynamic scene experiments, the proposed method achieves the best PSNR on Robot, Coral, and Streaks, with an average PSNR of 30.54 dB, outperforming UDR-GS (30.05 dB) and 4DGS (29.63 dB). On the SeaThru-NeRF dataset, the proposed method is generally comparable to or better than the 4DGS baseline. No-reference image quality assessment (NIQE) results indicate that the proposed method achieves optimal statistical naturalness on dewatered reconstructions in most scenes, further validating the effectiveness of the underwater physics model and depth supervision. Ablation studies reveal that the joint introduction of underwater physical modeling and monocular depth supervision does not consistently yield positive gains across all scenes; optimization interference exists between the two, and the full configuration is not always optimal.
\end{spacing}
\end{abstract}


\begin{keywords}
3D Gaussian splatting; underwater scene reconstruction; dynamic scene modeling; underwater physical imaging; depth supervision
\end{keywords}
